\chapter{Linear Algebra}

\section{Matrix-vector multiplication}

A $2\times 2$ matrix $A=\begin{pmatrix}a&b\\c&d\end{pmatrix}$ acts on
a vector $v=\begin{pmatrix}x\\y\end{pmatrix}$ by
\[
Av = x\begin{pmatrix}a\\c\end{pmatrix} + y\begin{pmatrix}b\\d\end{pmatrix}
= \begin{pmatrix}ax+by\\cx+dy\end{pmatrix}.
\]
That is, $Av$ is the combination of $A$'s two columns, $\begin{pmatrix}a\\c\end{pmatrix}$ and $\begin{pmatrix}b\\d\end{pmatrix}$, weighted by $v$'s entries $x,y$.

\section{The determinant as an area}

Take two vectors in the plane, $u=(a,c)$ and $v=(b,d)$ --- the columns
of $A$ from Section 1. They span a parallelogram.

\begin{center}
\begin{tikzpicture}[scale=0.9]
  \draw[->] (0,0) -- (4.5,0);
  \draw[->] (0,0) -- (0,3.5);
  \coordinate (O) at (0,0);
  \coordinate (U) at (3,1);
  \coordinate (V) at (1,2.5);
  \coordinate (W) at ($(U)+(V)$);
  \draw[thick,-{Stealth[]}] (O) -- (U) node[below right] {$u=(a,c)$};
  \draw[thick,-{Stealth[]}] (O) -- (V) node[above left] {$v=(b,d)$};
  \draw[dashed] (U) -- (W) -- (V);
  \draw[fill=gray!20,opacity=0.5] (O)--(U)--(W)--(V)--cycle;
  \draw[dotted] (0,0) rectangle (4,3.5);
\end{tikzpicture}
\end{center}

Enclose the parallelogram in the axis-aligned bounding box of width
$a+b$ and height $c+d$. The box consists of the parallelogram, two
right triangles with legs $a,c$ (combined area $ac$), two right
triangles with legs $b,d$ (combined area $bd$), and two $b\times c$
rectangles (combined area $2bc$). So
\[
(a+b)(c+d) = \text{Area} + ac + bd + 2bc
\quad\Longrightarrow\quad
\text{Area} = ad - bc.
\]

\begin{definition}
The \emph{determinant} of a matrix
$A=\begin{pmatrix}a&b\\c&d\end{pmatrix}$ is the \emph{signed} area
spanned by $u=(a,c)$, $v=(b,d)$. 
\end{definition}
We had just saw that $\det A = ad-bc$.
Geometrically: $\det A = 0$ exactly when the columns of $A$ are
parallel --- the parallelogram they span is flat, with zero area. 

\section{Eigenvalues of a $2\times 2$ matrix}

\begin{definition}
A nonzero vector $v$ is an \emph{eigenvector} of $A$ with
\emph{eigenvalue} $\lambda$ if $Av = \lambda v$.
\end{definition}

Rearranging, $(A-\lambda I)v = 0$ for some $v\ne 0$. We saw that
$(A-\lambda I)v$ is the combination of $(A-\lambda I)$'s two columns
weighted by $v$'s entries; for a \emph{nonzero} $v$ to give $0$, that
combination must be nontrivial, which is only possible if the two
columns are parallel, i.e.
\[
\det(A-\lambda I) = 0.
\]

\paragraph{Useful fact.}
The determinant is the product of eigenvalues, whereas the trace is
their sum.

\section{Eigendecomposition}
Suppose $A$ has eigenvectors $v_1,v_2$ with eigenvalues
$\lambda_1,\lambda_2$, and let $P=\begin{pmatrix}v_1 & v_2\end{pmatrix}$
(eigenvectors as columns). Then
\[
AP = \begin{pmatrix}Av_1 & Av_2\end{pmatrix}
   = \begin{pmatrix}\lambda_1 v_1 & \lambda_2 v_2\end{pmatrix}
   = PD, \qquad D=\begin{pmatrix}\lambda_1&0\\0&\lambda_2\end{pmatrix},
\]
so, if $v_1,v_2$ are independent ($P$ invertible),
\[
A = PDP^{-1}.
\]
$P$ is a \emph{change of basis}: $P^{-1}$ converts a vector from the
standard basis to \emph{eigenbasis} coordinates (coefficients along
$v_1,v_2$), and $P$ converts back. In eigenbasis coordinates $A$ acts
as the diagonal $D$.

\section{The determinant for a $3\times 3$ matrix}
\[
\det\begin{pmatrix}a_1&b_1&c_1\\a_2&b_2&c_2\\a_3&b_3&c_3\end{pmatrix}
= a_1\det\begin{pmatrix}b_2&c_2\\b_3&c_3\end{pmatrix}
- a_2\det\begin{pmatrix}b_1&c_1\\b_3&c_3\end{pmatrix}
+ a_3\det\begin{pmatrix}b_1&c_1\\b_2&c_2\end{pmatrix}.
\]
The same interpretations hold: the determinant is (signed) volume, the
determinant is the product of eigenvalues, and the trace is the sum of
eigenvalues.
